\backhead{FURTHER EXERCISES}{khred}
\begin{exlist}
\item A plant crown has semi-axes 0.35, 0.30 and 0.45 m. Find the ellipsoid volume. After drought the horizontal axes fall by 20 per cent: find the new volume and the fall.\exans{0.198 m$^3$; 0.127 m$^3$; 36 per cent}
\item A ring of 8 cameras surrounds a plant. Find the angle between neighbouring views. If each camera sees $\pm30^\circ$ about its axis, how many cameras see each point of the outer crown?\exans{45$^\circ$; one or two (overlap $15^\circ$)}
\item A leaf rolls when its water potential falls below $-1.2$ MPa. Convert this to a suction in cm of water and in pF.\exans{12\,232 cm; pF 4.09}
\item The deterministic plant gains 3.5 g of dry matter per day when well watered and 2.1 g under drought. Over 140 days, find the final masses if both start at 0.5 g and the drought begins on day 60.\exans{490.5 g and 378.5 g}
\item A scan has a point spacing of 2 mm. How many points cover a leaf of 300 cm$^2$?\exans{7500}
\end{exlist}
\backhead{ESSAY QUESTIONS}{khred}
{\small\itshape Answer each question in about one page, with equations, a sketch and at least one number from this chapter.}\par\vspace{3pt}
\begin{exlist}
\item Explain what phenotyping is and why a virtual, deterministic phenotyping experiment is useful for teaching.
\item Describe the steps of a three-dimensional scan, from the views to the measured traits.
\item Discuss how drought appears in the traits of a plant and how a scanner can detect it before yield falls.
\item Compare the deterministic and the volumetric phenotyping layers: what each computes and how they are linked.
\item Discuss the limits of a virtual scanner compared with a real one.
\item How could phenotyping data be used to set an irrigation trigger?
\end{exlist}
